% Build twice: pdflatex -interaction=nonstopmode -halt-on-error
%   -no-shell-escape INDEPENDENT_CHECKER_REPORT.tex
\documentclass[10pt,a4paper]{article}
\usepackage[T1]{fontenc}
\usepackage[utf8]{inputenc}
\usepackage{lmodern}
\usepackage[margin=2.3cm]{geometry}
\usepackage{amsmath,amssymb}
\usepackage[hidelinks]{hyperref}
\setlength{\parindent}{0pt}
\setlength{\parskip}{0.5\baselineskip}
\emergencystretch=3em
\sloppy
\urlstyle{tt}
\newcommand{\src}[1]{\href{#1}{\nolinkurl{#1}}}
\newcommand{\code}[1]{\texttt{\detokenize{#1}}}
\title{The independent F* proof checker:\\guarantee and interpretation architecture}
\author{}
\date{15 November 2026}
\begin{document}
\maketitle

\begin{abstract}
Core and Definition are two routes through the same independent, sort-indexed
F* proof checker.  If either certifies a proof document, its original
assertions cannot all be true in any interpretation satisfying the root-model
contract.  This report describes that precise theorem, the construction that
makes the model contract inhabited, the dependent value representation, replay
and recovery, and the boundary between the verified document theorem and the
native solver integration.  It describes the development on this PR, not a
claim that the migration or benchmark parity is complete.
\end{abstract}

\section{The guarantee}
Let \(\Sigma\) be a checked datatype schema, \(D\) a proof document, \(A_D\)
its original assertions, \(u\) a family of nonempty uninterpreted carriers,
and \(M\) an interpretation over the corresponding sort-indexed values.  The
interfaces \src{Z3IndependentChecker.fsti} expose the following two theorems:
\begin{align*}
 &\operatorname{check}(\Sigma,D)=\mathsf{RootCertified}
  \;\land\; \operatorname{root\_model}(u,\Sigma,M)
  \;\Longrightarrow\; \neg\operatorname{root\_truths}(u,\Sigma,M,A_D),\\
 &\operatorname{check\_definition}(\Sigma,D)=\mathsf{RootCertified}
  \;\land\; \operatorname{root\_model}(u,\Sigma,M)
  \;\Longrightarrow\; \neg\operatorname{root\_truths}(u,\Sigma,M,A_D).
\end{align*}
The first is \code{lemma_certified}; the second is
\code{lemma_certified_definition}.  Both quantify over the same \(u\),
\(\Sigma\), \(D\), and \emph{one interpretation for the whole document}.
\(\operatorname{root\_truths}\) is the conjunction of root truth for the
original assertions, not a statement merely about auxiliary replay formulas.
Thus no model in the stated class satisfies all those assertions.  For a
covered, well-typed document this is the intended refutation guarantee relative
to that model class; it is not an unrestricted theorem about arbitrary SMT-LIB
parsing, native code, or every possible theory.

The model hypothesis is not just an assumed, potentially empty contract.\linebreak[2]
\code{Z3TypedRoot.lemma_standard_root} constructs a root model
\(\operatorname{standard}(u,\Sigma,F)\) for every checked schema, admissible
carrier family \(u\), and choice of free symbols \(F\).  F* proves that this
model satisfies the root laws.  This is an inhabitation theorem; it does not
by itself identify all unsupported symbols with their full SMT-LIB theory.
See \src{Z3TypedRoot.fst} and \src{Z3SortedRootDocument.fst}.

\paragraph{Verdicts.}
The OCaml driver maps \code{RootCertified} to \code{ACCEPT},
\code{RootPending n} to \code{PARTIAL n}, \code{RootUncovered} to
\code{UNKNOWN} \code{unsupported-root-fragment}, and \code{RootRejected} to
\code{REJECT} (\src{ocaml_runtime/Z3IndependentMain.ml}).  Only
\code{ACCEPT} invokes the certification guarantee above.  A partial result
reports unresolved obligations, not a proof of satisfiability or a
counterexample.  An uncovered result is not certification, even if a lower
replay check found no pending obligations.  Solver \code{unknown}, timeouts,
errors, crashes, and unrun cases are not accepted proofs.

\section{From proof to verdict}
\begin{center}
\fbox{\parbox{0.90\linewidth}{\centering
Z3 assertions and candidate proof
\(\longrightarrow\) native marshaling and typed schema/document
\(\longrightarrow\) verified document replay and coverage gate
\(\longrightarrow\) Core verdict
\(\longrightarrow\) optional Definition recovery}}
\end{center}
The solver is a proof producer, not an oracle whose answer is accepted on
faith.  The document records original assertions, proof entries, and a root.
The checker validates declarations and assertions, replays entries using
earlier checked facts, and checks that the root is closed \code{false}.
Unverified or undecided steps leave explicit obligations.  The document
diagnostic records either pending obligations or a rejection stage
(\src{Z3DocumentSyntax.fst}, \src{Z3SortedProof.fst}).

Core is \code{Z3IndependentChecker.check}, the \code{sorted-core} route.
It uses the root replay and its verified certificate fallback.  Definition is
\code{check_definition}, selected by \code{--definition-replay}; it first
runs Core and attempts additional verified recovery only if Core is pending.
It retains any Core certified, uncovered, or rejected result.  Successful
recovery certifies the original assertions under the same general root-model
contract, rather than silently changing the meaning of \code{ACCEPT}.
The executable dispatch and exact reference behavior are in
\src{ocaml_runtime/Z3IndependentMain.ml} and
\src{Z3IndependentChecker.fsti}.

The root gate also checks its supported schema, scopes and document-term
coverage.  In the underlying root document route, a structurally valid closed
proof with no pending obligations is still \code{RootUncovered} if that
coverage test fails.  See \code{root_covered} and
\code{check_root_document} in \src{Z3SortedRootDocument.fst}.  The more
general Core and Definition entry points use their own verified recovery
contracts, so their certification should be read through the exported
theorems, not identified with a single internal check.

\section{Sort-indexed values and evaluation}
\subsection{The carrier family}
The value infrastructure is parameterized by a dependent family
\[
  \mathit{base} : \mathit{sort}\to\mathit{Type0}.
\]
This maps each sort \emph{AST} to an F* type.  It is not a single untyped
sum of all SMT values, and Definition is not a different value domain from
Core.  \src{Z3DatatypeValue.fst} represents values as trees:
\[
  \mathit{tree}(\mathit{base})
    = \mathsf{Scalar}(s,\mathit{base}(s))
      \mid \mathsf{Constructor}(f,\mathit{children}).
\]
The type \code{value base sigma s} refines a tree by the executable
\code{has_sort} predicate.  A scalar must carry the requested non-datatype
sort; a constructor must occur in the datatype signature \(\sigma\), have
the requested range, and have recursively well-sorted fields.  Argument
lists are similarly refined against a list of domain sorts.  Constructor
disjointness, injectivity, and the matched-selector laws are proved over
these trees; they are not consequences of comparing printed names.

For the root model, \code{Z3TypedRoot.carrier u s} instantiates the family:
Boolean, integer, and real sorts use F* \code{bool}, \code{int}, and
\code{real}; valid bit-vector and finite-field sorts use indexed carriers;
sequences use lists; and datatype-free array sorts use spaces of ghost
functions over typed index tuples.  Other sorts use a nonempty
\code{u.sort_carrier s}.  Datatype values are checked constructor trees,
not scalar payloads at the datatype sort.  A concrete default inhabitant is
constructed recursively.  The existence of a carrier does not itself imply
that every operator on it has been given full theory semantics.

\subsection{Interpretations and models}
\code{Z3DatatypeValue.interpretation base sigma} assigns a typed result to
each function application and a typed result to each quantified or lambda
body.  A scoped environment maps each bound-variable index to a value of
the corresponding sort.  Ghost evaluation recursively evaluates a
well-scoped expression to a value at \emph{its} AST sort.  Registered
constructors build trees; other applications and binders use the shared
interpretation.  These semantic definitions are ghost specifications, not
runtime procedures for deciding equality of arbitrary real or function values.

\code{Z3TypedRoot.root_model} requires one interpretation to satisfy its
quantifier, interpreted-operator, lambda, constructor, recognizer, and
matched-accessor laws.  The standard interpretation composes the standard
theory and datatype interpretations with free symbol choices.  For example,
it fixes the designated Boolean, integer and real operations, and gives
datatype-free arrays functional select/store/lambda semantics where their
law preconditions apply.  Symbols not fixed by the contract remain
model choices; selectors applied to the wrong constructor are an example.
Theory coverage is therefore a property of the actual laws and guards, not
of a sort name alone.  The typed definitions and their law preconditions
are in \src{Z3TypedRoot.fst}, \src{Z3SortedSemantics.fst}, and the
associated sorted-theory modules.

\section{Definitions, binders, and cached replay}
Definition recovery recognizes eligible definition introductions and can
construct augmented documents for definitions, NNF, Skolemization, dependent
contexts, or selected binder replay.  Eligibility, freshness, scope, and
root certification checks gate the candidates.  In particular, fresh
definitions and Skolem choices require a \emph{model-extension} argument:
the extended interpretation satisfies the auxiliary material while
preserving truth of the original assertions.  It is not enough to assume
that every introduced symbol already has the desired interpretation.
\src{Z3IndependentDefinitionModel.fsti} specifies an override that
preserves other applications, quantification, and fresh input expressions;
the document-level proof combines the relevant preservation lemmas with
replay soundness.  See also \src{Z3IndependentSkolemModel.fst} and
\src{Z3IndependentDefinitionSkolemDocument.fsti}.

The cached replay is an optimization, not an additional semantic axiom.
\src{Z3IndependentCachedReplay.fsti} states exact equality of its
document diagnostic with \code{Z3SortedProof.diagnose_document}, and
exact equality of \code{check_document} with
\code{Z3SortedProof.check_document}.  The cache is indexed by the schema
and original assertion list.  The Definition and Skolem paths may use
this replay, but their public contracts retain the reference verdict and
soundness theorem.

\section{Limits and trust boundary}
\begin{itemize}
\item The theorem is conditional on a checked schema and
\code{root_model}; the standard-model theorem establishes an inhabitant,
but the claimed standard meaning is restricted to supported schemas,
expressions, and the stated theory laws.  In particular, nested arrays
mentioning registered datatypes and zero-index arrays are excluded by
the supported-sort gate.  Opaque operations are not automatically given
their full SMT-LIB semantics.
\item The separate experimental Definition-plus-array entry point has a
different theorem: certification is shown for \code{TR.standard} models,
not for every arbitrary \code{root_model}.  It must not be presented as
the generic Definition theorem or silently substituted for that route.
\item The proof is about the F* schema and document actually checked.
The native Z3 API, foreign bindings and marshaling, assertion/proof
transport, extraction, OCaml compiler and runtime, F* verification
toolchain and its SMT backend, and the execution environment remain part
of the end-to-end trust boundary.  Correct transport from the original
SMT-LIB session to the checked document is not established by the two
interface theorems alone.
\item Soundness is not completeness or a performance guarantee.  Pending
rules, unsupported terms, solver differences, and resource limits can
prevent certification.  The independent migration, retirement of
\code{Z3ValueTheory.fst}, and full-corpus legacy parity remain work in
progress; see \src{PORTING_PROGRESS.md}.
\end{itemize}

\section{Reading the implementation}
The contracts to start from are \src{Z3IndependentChecker.fsti}, the
root gate and theorem in \src{Z3SortedRootDocument.fst}, the typed
carrier and root laws in \src{Z3TypedRoot.fst}, and the value/evaluator
in \src{Z3DatatypeValue.fst}.  The root-model statements are semantic
F* lemmas, while the proof-document decisions and coverage checks are
executable.  Build and benchmark status, including what was actually
verified and measured, is tracked separately in \src{PORTING_PROGRESS.md}.

\end{document}
